% !TeX root = ../main.tex
\section{Zero-Queue Baselines}\label{sec:zero-queue}

\begin{hypothesis}
  At QD=1 a single I/O is in flight, so latency is the full host-stack + controller + flash path. We
  expect \qty{4}{\kibi\byte} random reads in the tens of microseconds (flash read dominates) and
  \qty{128}{\kibi\byte} sequential transfers to be bandwidth-limited at a few hundred
  \unit{\micro\second}. Writes should be \emph{faster} than reads at QD=1 because the SSD's DRAM/SLC
  write buffer acknowledges them before the flash program completes.
\end{hypothesis}

\begin{knobs}
  Fixed       & iodepth 1; one job; \code{O_DIRECT} \\
  Cases       & 4~KiB random read, 4~KiB random write, 128~KiB sequential read, 128~KiB sequential write \\
  Script      & \code{SSD_Performance_Profiling/scripts/run.sh zero_queue} \\
\end{knobs}

\begin{table}[tbp]
  \centering\small
  \caption{QD=1 performance on the Micron 2210 NVMe SSD (medians of 3 reps; latencies in
  \si{\micro\second}). Random reads pay the full flash-read latency; random writes are faster because the
  write buffer absorbs them.}\label{tab:zero-queue}
  \begin{tabularx}{\linewidth}{@{}llrrrrrrX@{}}
    \toprule
    Pattern & Block & IOPS & MiB/s & Mean & p50 & p95 & p99 & p99.9 \\
    \midrule
    rand read  & \qty{4}{\kibi\byte}   & \num{20150} & \num{78.7}  & \num{49.6}  & \num{47.5} & \num{50.0} & \num{68}  & \num{460} \\
    rand write & \qty{4}{\kibi\byte}   & \num{68037} & \num{265.8} & \num{14.7}  & \num{12.5} & \num{14.3} & \num{26}  & \num{440} \\
    seq read   & \qty{128}{\kibi\byte} & \num{1813}  & \num{226.6} & \num{551.4} & \num{549}  & \num{847}  & \num{1000}& \num{1349} \\
    seq write  & \qty{128}{\kibi\byte} & \num{5959}  & \num{744.9} & \num{167.8} & \num{94.2} & \num{396}  & \num{1417}& \num{2097} \\
    \bottomrule
  \end{tabularx}
\end{table}

\subsection{Analysis}
The QD=1 numbers confirm an NVMe SSD, not a rotational disk: \qty{4}{\kibi\byte} random read latency is
\qty{49.6}{\micro\second} (not the milliseconds a HDD would show), dominated by the NAND read plus the
NVMe/\code{libaio} host path rather than any seek. Random \emph{writes} are over three times faster
(\qty{14.7}{\micro\second}, \num{68}k IOPS) because the controller acknowledges them from its DRAM/SLC
write buffer before the flash program finishes---a latency the QD=1 read cannot hide. The
\qty{128}{\kibi\byte} sequential rows are bandwidth-bound: at \qty{227}{\mebi\byte\per\second} (read) and
\qty{745}{\mebi\byte\per\second} (write) each transfer takes hundreds of microseconds, and the single
outstanding request leaves the device under-fed---\cref{sec:block-size,sec:queue-depth} show throughput
climbing sharply once larger blocks or deeper queues keep more data moving. Two caveats: QD=1 reads
cannot be prefetched by the device for a cold random stream, so this is a genuine floor; and the p99.9
column already shows occasional \qty{440}{}--\qty{2100}{\micro\second} outliers from background flush/GC,
previewing the tail behaviour in \cref{sec:tail}.
